\section{Data}
\label{sec:data}

\subsection{Data Sourcing}

Following Alvi's design (\S3.1 of \autocite{alvi2018pgm}), the candidate variable set is split into three groups:

\begin{itemize}
  \item \textbf{Financial (5 series).} Daily closes from Yahoo Finance: WTI and Brent continuous front-month futures, the ICE
    US Dollar Index, the S\&P 500, and the 10-year US Treasury yield.
  \item \textbf{Macroeconomic (8 series).} Monthly series from the Federal
    Reserve Economic Data (FRED) API, matching Alvi's Table 3.2 exactly:
    energy CPI, oil \& gas extraction capacity and capacity utilisation,
    three industrial-production indices (oil \& gas extraction, crude
    petroleum \& natural gas, and oil \& gas support activities), total
    industrial production, and the oil \& gas extraction producer price
    index.
  \item \textbf{Geopolitical / physical-market (11 series).} Monthly Short-Term
    Energy Outlook (STEO) series from the EIA Open Data v2 API, matching 11
    of the 14 series in Alvi's Table 3.1 (OPEC/non-OPEC supply, OECD/non-OECD
    demand, US production and refinery input, OPEC spare/production
    capacity, and OECD/non-OECD stock changes). The remaining 3 of Alvi's 14
    series (\texttt{STEO.RGDPQ\_NONOECD.M}, \texttt{STEO.RGDPQ\_OECD.M},
    \texttt{STEO.FOREX\_WORLD.M}) return no data against the current STEO
    vintage and have been retired or renamed since 2018; rather than guess a
    replacement series ID, we omit them and document the omission here
    explicitly.
\end{itemize}

One data-quality nuance specific to the EIA STEO series is worth flagging explicitly: each pull includes roughly a year of the EIA's own forward-looking forecast beyond the present day, which was truncated before any exploratory or stylised-facts analysis so that only realised history is analysed. The resulting data dictionary is given in Table~\ref{tab:data-dictionary}.

\begingroup
\small
\setlength{\tabcolsep}{3pt}
\begin{longtable}{>{\raggedright\arraybackslash}p{1.9cm}>{\raggedright\arraybackslash}p{2.3cm}>{\raggedright\arraybackslash}p{4.2cm}>{\raggedright\arraybackslash}p{1.4cm}>{\raggedright\arraybackslash}p{0.8cm}>{\raggedright\arraybackslash}p{1.45cm}>{\raggedright\arraybackslash}p{1.45cm}}
\caption{GWP1 data dictionary: 24 series across the financial, macroeconomic, and geopolitical/physical-market groups.} \label{tab:data-dictionary} \\
\toprule
Name & Category & Source & Freq. & Units & Start & End \\
\midrule
\endfirsthead
\caption[]{GWP1 data dictionary: 24 series across the financial, macroeconomic, and geopolitical/physical-market groups.} \\
\toprule
Name & Category & Source & Freq. & Units & Start & End \\
\midrule
\endhead
\midrule
\multicolumn{7}{r}{Continued on next page} \\
\midrule
\endfoot
\bottomrule
\endlastfoot
brent\_\allowbreak futures & financial & Yahoo Finance & daily & n/a & 2007-07-30 & 2026-10-06 \\
sp500 & financial & Yahoo Finance & daily & n/a & 1990-01-02 & 2026-10-05 \\
usd\_\allowbreak index & financial & Yahoo Finance & daily & n/a & 1990-01-01 & 2026-10-06 \\
ust10y\_\allowbreak yield & financial & Yahoo Finance & daily & n/a & 1990-01-02 & 2026-10-05 \\
wti\_\allowbreak futures & financial & Yahoo Finance & daily & n/a & 2000-08-23 & 2026-10-06 \\
demand\_\allowbreak non\_\allowbreak oecd & geopolitical & U.S. Energy Information Administration (EIA) & monthly & n/a & 1990-01-01 & 2026-10-01 \\
demand\_\allowbreak oecd & geopolitical & U.S. Energy Information Administration (EIA) & monthly & n/a & 1990-01-01 & 2026-10-01 \\
non\_\allowbreak oecd\_\allowbreak stock\_\allowbreak change & geopolitical & U.S. Energy Information Administration (EIA) & monthly & n/a & 2003-02-01 & 2026-10-01 \\
oecd\_\allowbreak inventory\_\allowbreak balance & geopolitical & U.S. Energy Information Administration (EIA) & monthly & n/a & 2003-01-01 & 2026-10-01 \\
oecd\_\allowbreak stock\_\allowbreak change & geopolitical & U.S. Energy Information Administration (EIA) & monthly & n/a & 2003-02-01 & 2026-10-01 \\
opec\_\allowbreak production\_\allowbreak capacity & geopolitical & U.S. Energy Information Administration (EIA) & monthly & n/a & 1990-01-01 & 2026-10-01 \\
opec\_\allowbreak spare\_\allowbreak capacity & geopolitical & U.S. Energy Information Administration (EIA) & monthly & n/a & 1993-01-01 & 2026-10-01 \\
supply\_\allowbreak non\_\allowbreak opec & geopolitical & U.S. Energy Information Administration (EIA) & monthly & n/a & 1993-01-01 & 2026-10-01 \\
supply\_\allowbreak opec & geopolitical & U.S. Energy Information Administration (EIA) & monthly & n/a & 1993-01-01 & 2026-10-01 \\
us\_\allowbreak crude\_\allowbreak production & geopolitical & U.S. Energy Information Administration (EIA) & monthly & n/a & 1990-01-01 & 2026-10-01 \\
us\_\allowbreak refinery\_\allowbreak input & geopolitical & U.S. Energy Information Administration (EIA) & monthly & n/a & 1990-01-01 & 2026-10-01 \\
capacity\_\allowbreak oil\_\allowbreak gas\_\allowbreak extraction & macroeconomic & FRED (Federal Reserve Economic Data) & monthly & n/a & 1972-01-01 & 2026-08-01 \\
capacity\_\allowbreak util\_\allowbreak oil\_\allowbreak gas\_\allowbreak extraction & macroeconomic & FRED (Federal Reserve Economic Data) & monthly & n/a & 1972-01-01 & 2026-08-01 \\
cpi\_\allowbreak energy & macroeconomic & FRED (Federal Reserve Economic Data) & monthly & n/a & 1957-01-01 & 2026-08-01 \\
indprod\_\allowbreak crude\_\allowbreak petroleum\_\allowbreak ng & macroeconomic & FRED (Federal Reserve Economic Data) & monthly & n/a & 1972-01-01 & 2021-04-01 \\
indprod\_\allowbreak oil\_\allowbreak gas\_\allowbreak extraction & macroeconomic & FRED (Federal Reserve Economic Data) & monthly & n/a & 1972-01-01 & 2026-08-01 \\
indprod\_\allowbreak oil\_\allowbreak gas\_\allowbreak support & macroeconomic & FRED (Federal Reserve Economic Data) & monthly & n/a & 1972-01-01 & 2026-08-01 \\
indprod\_\allowbreak total & macroeconomic & FRED (Federal Reserve Economic Data) & monthly & n/a & 1919-01-01 & 2026-08-01 \\
ppi\_\allowbreak oil\_\allowbreak gas\_\allowbreak extraction & macroeconomic & FRED (Federal Reserve Economic Data) & monthly & n/a & 1985-12-01 & 2026-08-01 \\
\end{longtable}

\endgroup

\subsection{Data Cleaning}
\label{sec:data-cleaning}

Every series is passed through the same three-stage:

\begin{enumerate}
  \item \textbf{Bad data.} Non-monotonic or duplicated timestamps, and
    economically impossible non-positive values, are treated as errors and
    set to missing. One deliberate exception is built in: the OECD and
    non-OECD inventory \emph{stock-change} series are genuinely
    \emph{signed} flow variables (a negative value denotes a net stock
    \emph{build}, not an error), so the non-positive check is disabled for
    those two series only.
  \item \textbf{Extreme outliers.} A robust, median-absolute-deviation
    (MAD) based modified $z$-score flags (but does not remove) observations
    with $|z| > 3.5$, so genuine price shocks remain visible in the
    sterilised series while still being documented.
  \item \textbf{Missing values.} Slow-moving macroeconomic/geopolitical
    series are imputed by linear interpolation (appropriate given their
    short, infrequent gaps); the daily financial series instead use a
    volatility-matched Brownian-bridge simulation between the surrounding
    observed values, which avoids injecting artificial flatness into a price
    series across a longer gap (e.g.\ a data-vendor outage spanning several
    trading days).
\end{enumerate}

One specific finding is worth flagging on its own merits rather than leaving implicit in the narrative above: the raw WTI futures feed contains exactly one non-positive observation, corresponding to the historic 20 April 2020
episode when the expiring WTI contract briefly traded at a negative price as storage capacity was exhausted and long holders could not take delivery. Our general-purpose bad-data rule (non-positive price is an error) is
a defensible default for a panel of 24 heterogeneous series, but it has the side effect of treating this genuine, famous market event identically to a data error and imputing it away. We flag this transparently here, and again in Section~\ref{sec:stylized-facts}, rather than silently losing the observation.
